\documentclass[11pt,a4paper]{article}
\usepackage[utf8]{inputenc}
\usepackage[T1]{fontenc}
\usepackage{lmodern}
\usepackage[french]{babel}
\usepackage{amsmath,amssymb,mathtools}
\usepackage{icomma}
\usepackage{xcolor}
\usepackage[most]{tcolorbox}
\usepackage{enumitem}
\usepackage[a4paper,margin=2.2cm,headheight=26pt,headsep=10pt]{geometry}
\usepackage{fancyhdr}
\usepackage[hidelinks]{hyperref}
\definecolor{primary}{RGB}{222,49,99}
\definecolor{primarydark}{RGB}{150,22,55}
\tcbset{boxbase/.style={enhanced,breakable,boxrule=0.9pt,arc=2.5pt,
  left=3mm,right=3mm,top=2.3mm,bottom=2.3mm,fonttitle=\bfseries,coltitle=white}}
\newtcolorbox[auto counter]{exo}[1][]{boxbase,colback=primary!4,colframe=primary,
  title={\textsc{Exercice}~\thetcbcounter\ifx\relax#1\relax\relax\else\textmd{ --- \itshape #1}\fi}}
\newcommand{\R}{\mathbb{R}}
\newcommand{\ds}{\displaystyle}
\pagestyle{fancy}\fancyhf{}
\fancyhead[L]{\small\textcolor{primarydark}{\textbf{Thème 3} — Factorisation et résolution}}
\fancyhead[R]{\small\textcolor{primarydark}{\textsc{Difficile}}}
\fancyfoot[C]{\small\thepage}
\renewcommand{\headrulewidth}{0.8pt}
\begin{document}

\begin{center}
{\LARGE\bfseries\color{primarydark} Niveau Difficile — Exercices}\\[2pt]
{\color{primary}\rule{0.45\textwidth}{1.2pt}}\\[3pt]
{\small\itshape Niveau concours. Les sous-questions construisent la difficulté pas à pas.}
\end{center}
\vspace{1mm}

\begin{exo}[factorisation complète et fraction rationnelle]
Soit $P(x)=4x^{3}-8x^{2}+5x-1$.
\begin{enumerate}[label=\textbf{\alph*)},leftmargin=2.2em,topsep=1pt,itemsep=2pt]
  \item Montre que $x=1$ est racine, puis utilise le schéma de Horner pour écrire $P(x)=(x-1)(4x^{2}-4x+1)$.
  \item Reconnais un carré parfait dans $4x^{2}-4x+1$ et donne la factorisation \emph{complète} de $P(x)$.
  \item Résous $P(x)=0$ dans $\R$.
  \item Simplifie $\ds\frac{P(x)}{2x^{2}-3x+1}$ en précisant la condition d'existence.
\end{enumerate}
\end{exo}

\begin{exo}[reconstituer un trinôme à partir de ses racines]
On considère $P(x)=ax^{2}+bx+c$ avec $a\neq 0$. On sait que $x=1$ et $x=3$ sont les deux racines de $P$,
et que $P(2)=2$.
\begin{enumerate}[label=\textbf{\alph*)},leftmargin=2.2em,topsep=1pt,itemsep=2pt]
  \item Écris $P(x)$ sous forme factorisée en fonction de $a$.
  \item Utilise $P(2)=2$ pour déterminer $a$.
  \item Développe $P$, donne $b$ et $c$, puis calcule $P(0)$.
\end{enumerate}
\end{exo}

\begin{exo}[discriminant et paramètre]
Soit $m$ un réel. On considère l'équation $\ (m-3)x^{2}-2x+2=0$.
\begin{enumerate}[label=\textbf{\alph*)},leftmargin=2.2em,topsep=1pt,itemsep=2pt]
  \item Que devient l'équation lorsque $m=3$~? Combien admet-elle alors de solution(s)~?
  \item Pour $m\neq 3$, exprime le discriminant $\Delta$ en fonction de $m$.
  \item Détermine l'ensemble des valeurs de $m$ pour lesquelles l'équation admet \emph{exactement deux racines réelles distinctes}. (Attention à bien exclure le cas $m=3$.)
\end{enumerate}
\end{exo}

\end{document}
